% --- Archon physics formalization source begin ---
% archon:physics
% archon:covers PhyXMiniProblems/problem_phyx_mini_0274.lean
% archon:source-report reports/phyx_mini/problem_phyx_mini_0274.source.json

\chapter{Physics problem phyx\_mini\_0274}
\label{ch:PhyXMiniProblems_problem_phyx_mini_0274}

\paragraph{Problem source.}
\#\# Physical scenario

Figure shows the kinetic energy \$K\$ of a simple pendulum versus its angle \$\textbackslash{}theta\$ from the vertical. The vertical axis scale is set by \$K\_s = 10.0\textbackslash{} mJ\$. The pendulum bob has mass \$0.200\textbackslash{} kg\$.

\#\# Question

What is the length of the pendulum?

\#\# Answer choices

- A: A: 1.47 m

- B: B: 1.50 m

- C: C: 1.53 m

- D: D: 1.56 m

\#\# Figure caption (auxiliary; use the image as primary evidence)

The image is a graph with a thick parabolic curve. Here's the detailed content:

- **Axes**: 
  - The horizontal axis is labeled as \textbackslash{}( \textbackslash{}theta \textbackslash{}) (mrad), ranging from -100 to 100, with tick marks at -100, -50, 0, 50, and 100.
  - The vertical axis is labeled as \textbackslash{}( K \textbackslash{}) (mJ).

- **Curve**: 
  - The curve is symmetrical and peaks at the center, aligned with the vertical axis.
  - The apex of the curve is labeled with the value \textbackslash{}( K\_s \textbackslash{}).

- **Grid**: 
  - A light grid is present in the background which aids in estimating the coordinates.

Overall, the graph represents a parabola indicating a relationship between \textbackslash{}( \textbackslash{}theta \textbackslash{}) (angle in milliradians) and \textbackslash{}( K \textbackslash{}) (energy in millijoules).

\#\# Dataset metadata

Resonance and Harmonics; Physical Model Grounding Reasoning, Multi-Formula Reasoning

\paragraph{Recorded answer/context.}
C: C: 1.53 m

\paragraph{Figure/image path.}
/root/proposal\_for\_physic/science-mango/phyx\_mini\_run/phyx\_data/test\_image/274.png

\paragraph{Formalization target.}
create a compiling Lean file with sorry bodies at `PhyXMiniProblems/problem\_phyx\_mini\_0274.lean`. The Lean declarations must preserve the physical quantities, dimensions or dimensional roles, figure labels, governing-law hypotheses, and final relation expressed by this problem.
Use Mathlib/Physlib names found through LeanExplore where available. If a physics API is missing, introduce faithful local abstractions rather than scalar placeholder aliases.

\begin{theorem}[Physics formalization target]
\label{thm:physics:phyx_mini_0274:target}
\lean{PhyXMiniProblems.ProblemPhyXMini0274.problem_phyx_mini_0274}
\uses{def:physics:phyx-mini-0274:phyxminiproblems-problemphyxmini0274-lengthinmeters, def:physics:phyx-mini-0274:phyxminiproblems-problemphyxmini0274-simplependulumsetup, def:physics:phyx-mini-0274:phyxminiproblems-problemphyxmini0274-matchesproblemandprimaryfigure, def:physics:phyx-mini-0274:phyxminiproblems-problemphyxmini0274-usesstandardearthgravity, def:physics:phyx-mini-0274:phyxminiproblems-problemphyxmini0274-hasphysicalpendulumparameters, def:physics:phyx-mini-0274:phyxminiproblems-problemphyxmini0274-satisfiesidealpendulumenergylaw, def:physics:phyx-mini-0274:phyxminiproblems-problemphyxmini0274-recordedanswerchoice, def:physics:phyx-mini-0274:phyxminiproblems-problemphyxmini0274-matchesdisplayedlength, def:physics:phyx-mini-0274:phyxminiproblems-problemphyxmini0274-isuniqueclosestdisplayedchoice}
The assigned autoformalize agent should translate this physics problem into Lean declarations in the covered file, with theorem and lemma proof bodies written as `by sorry`.
\end{theorem}
\begin{proof}
This is an autoformalization task, not a proof task. Produce statements that can later be proved by the physics prover without weakening the physical model.
\end{proof}

% --- Archon declaration topology begin ---
\section{Declaration topology}
\begin{definition}[Mass Quantity]
  \label{def:physics:phyx-mini-0274:phyxminiproblems-problemphyxmini0274-massquantity}
  \lean{PhyXMiniProblems.ProblemPhyXMini0274.MassQuantity}
  A nonnegative physical mass
\end{definition}
\begin{proof}
  This is the declared model definition of Mass Quantity.
\end{proof}
\begin{definition}[Length Quantity]
  \label{def:physics:phyx-mini-0274:phyxminiproblems-problemphyxmini0274-lengthquantity}
  \lean{PhyXMiniProblems.ProblemPhyXMini0274.LengthQuantity}
  A nonnegative physical length
\end{definition}
\begin{proof}
  This is the declared model definition of Length Quantity.
\end{proof}
\begin{definition}[Acceleration Quantity]
  \label{def:physics:phyx-mini-0274:phyxminiproblems-problemphyxmini0274-accelerationquantity}
  \lean{PhyXMiniProblems.ProblemPhyXMini0274.AccelerationQuantity}
  A nonnegative physical acceleration, with dimension L T⁻²
\end{definition}
\begin{proof}
  This is the declared model definition of Acceleration Quantity.
\end{proof}
\begin{definition}[Energy Quantity]
  \label{def:physics:phyx-mini-0274:phyxminiproblems-problemphyxmini0274-energyquantity}
  \lean{PhyXMiniProblems.ProblemPhyXMini0274.EnergyQuantity}
  Physical energy, with dimension M L² T⁻²
\end{definition}
\begin{proof}
  This is the declared model definition of Energy Quantity.
\end{proof}
\begin{definition}[mass In Kilograms]
  \label{def:physics:phyx-mini-0274:phyxminiproblems-problemphyxmini0274-massinkilograms}
  \lean{PhyXMiniProblems.ProblemPhyXMini0274.massInKilograms}
  \uses{def:physics:phyx-mini-0274:phyxminiproblems-problemphyxmini0274-massquantity}
  Kilogram readout of a physical mass
\end{definition}
\begin{proof}
  This is the declared model definition of mass In Kilograms.
\end{proof}
\begin{definition}[length In Meters]
  \label{def:physics:phyx-mini-0274:phyxminiproblems-problemphyxmini0274-lengthinmeters}
  \lean{PhyXMiniProblems.ProblemPhyXMini0274.lengthInMeters}
  \uses{def:physics:phyx-mini-0274:phyxminiproblems-problemphyxmini0274-lengthquantity}
  Metre readout of a physical length
\end{definition}
\begin{proof}
  This is the declared model definition of length In Meters.
\end{proof}
\begin{definition}[acceleration In Meters Per Second Squared]
  \label{def:physics:phyx-mini-0274:phyxminiproblems-problemphyxmini0274-accelerationinmeterspersecondsquared}
  \lean{PhyXMiniProblems.ProblemPhyXMini0274.accelerationInMetersPerSecondSquared}
  \uses{def:physics:phyx-mini-0274:phyxminiproblems-problemphyxmini0274-accelerationquantity}
  Metre-per-second-squared readout of a physical acceleration
\end{definition}
\begin{proof}
  This is the declared model definition of acceleration In Meters Per Second Squared.
\end{proof}
\begin{definition}[energy In Joules]
  \label{def:physics:phyx-mini-0274:phyxminiproblems-problemphyxmini0274-energyinjoules}
  \lean{PhyXMiniProblems.ProblemPhyXMini0274.energyInJoules}
  \uses{def:physics:phyx-mini-0274:phyxminiproblems-problemphyxmini0274-energyquantity}
  Joule readout of a physical energy
\end{definition}
\begin{proof}
  This is the declared model definition of energy In Joules.
\end{proof}
\begin{definition}[energy In Milli Joules]
  \label{def:physics:phyx-mini-0274:phyxminiproblems-problemphyxmini0274-energyinmillijoules}
  \lean{PhyXMiniProblems.ProblemPhyXMini0274.energyInMilliJoules}
  \uses{def:physics:phyx-mini-0274:phyxminiproblems-problemphyxmini0274-energyquantity, def:physics:phyx-mini-0274:phyxminiproblems-problemphyxmini0274-energyinjoules}
  Millijoule readout used on the graph's vertical axis
\end{definition}
\begin{proof}
  This is the declared model definition of energy In Milli Joules.
\end{proof}
\begin{definition}[milli Radians To Radians]
  \label{def:physics:phyx-mini-0274:phyxminiproblems-problemphyxmini0274-milliradianstoradians}
  \lean{PhyXMiniProblems.ProblemPhyXMini0274.milliRadiansToRadians}
  Convert a scalar graph coordinate in milliradians to radians
\end{definition}
\begin{proof}
  This is the declared model definition of milli Radians To Radians.
\end{proof}
\begin{definition}[Horizontal Axis Role]
  \label{def:physics:phyx-mini-0274:phyxminiproblems-problemphyxmini0274-horizontalaxisrole}
  \lean{PhyXMiniProblems.ProblemPhyXMini0274.HorizontalAxisRole}
  Physical quantity assigned to the graph's horizontal axis
\end{definition}
\begin{proof}
  This is the declared model definition of Horizontal Axis Role.
\end{proof}
\begin{definition}[Vertical Axis Role]
  \label{def:physics:phyx-mini-0274:phyxminiproblems-problemphyxmini0274-verticalaxisrole}
  \lean{PhyXMiniProblems.ProblemPhyXMini0274.VerticalAxisRole}
  Physical quantity assigned to the graph's vertical axis
\end{definition}
\begin{proof}
  This is the declared model definition of Vertical Axis Role.
\end{proof}
\begin{definition}[Angle Display Unit]
  \label{def:physics:phyx-mini-0274:phyxminiproblems-problemphyxmini0274-angledisplayunit}
  \lean{PhyXMiniProblems.ProblemPhyXMini0274.AngleDisplayUnit}
  Unit printed on the horizontal axis
\end{definition}
\begin{proof}
  This is the declared model definition of Angle Display Unit.
\end{proof}
\begin{definition}[Energy Display Unit]
  \label{def:physics:phyx-mini-0274:phyxminiproblems-problemphyxmini0274-energydisplayunit}
  \lean{PhyXMiniProblems.ProblemPhyXMini0274.EnergyDisplayUnit}
  Unit printed on the vertical axis
\end{definition}
\begin{proof}
  This is the declared model definition of Energy Display Unit.
\end{proof}
\begin{definition}[Graph Curve Shape]
  \label{def:physics:phyx-mini-0274:phyxminiproblems-problemphyxmini0274-graphcurveshape}
  \lean{PhyXMiniProblems.ProblemPhyXMini0274.GraphCurveShape}
  Qualitative shape of the heavy curve drawn in the supplied graph
\end{definition}
\begin{proof}
  This is the declared model definition of Graph Curve Shape.
\end{proof}
\begin{definition}[Pendulum Kinetic Energy Figure]
  \label{def:physics:phyx-mini-0274:phyxminiproblems-problemphyxmini0274-pendulumkineticenergyfigure}
  \lean{PhyXMiniProblems.ProblemPhyXMini0274.PendulumKineticEnergyFigure}
  \uses{def:physics:phyx-mini-0274:phyxminiproblems-problemphyxmini0274-energyquantity, def:physics:phyx-mini-0274:phyxminiproblems-problemphyxmini0274-horizontalaxisrole, def:physics:phyx-mini-0274:phyxminiproblems-problemphyxmini0274-verticalaxisrole, def:physics:phyx-mini-0274:phyxminiproblems-problemphyxmini0274-angledisplayunit, def:physics:phyx-mini-0274:phyxminiproblems-problemphyxmini0274-energydisplayunit, def:physics:phyx-mini-0274:phyxminiproblems-problemphyxmini0274-graphcurveshape}
  The labeled graph and the energy curve it displays. The angular argument of kineticEnergyAtRadians is an SI radian readout; the endpoint fields retain the milliradian coordinates printed in the image
\end{definition}
\begin{proof}
  This is the declared model definition of Pendulum Kinetic Energy Figure.
\end{proof}
\begin{definition}[Simple Pendulum Setup]
  \label{def:physics:phyx-mini-0274:phyxminiproblems-problemphyxmini0274-simplependulumsetup}
  \lean{PhyXMiniProblems.ProblemPhyXMini0274.SimplePendulumSetup}
  \uses{def:physics:phyx-mini-0274:phyxminiproblems-problemphyxmini0274-massquantity, def:physics:phyx-mini-0274:phyxminiproblems-problemphyxmini0274-lengthquantity, def:physics:phyx-mini-0274:phyxminiproblems-problemphyxmini0274-accelerationquantity, def:physics:phyx-mini-0274:phyxminiproblems-problemphyxmini0274-pendulumkineticenergyfigure}
  The simple pendulum and the graph supplied with the exercise
\end{definition}
\begin{proof}
  This is the declared model definition of Simple Pendulum Setup.
\end{proof}
\begin{definition}[Matches Problem And Primary Figure]
  \label{def:physics:phyx-mini-0274:phyxminiproblems-problemphyxmini0274-matchesproblemandprimaryfigure}
  \lean{PhyXMiniProblems.ProblemPhyXMini0274.MatchesProblemAndPrimaryFigure}
  \uses{def:physics:phyx-mini-0274:phyxminiproblems-problemphyxmini0274-massinkilograms, def:physics:phyx-mini-0274:phyxminiproblems-problemphyxmini0274-energyinjoules, def:physics:phyx-mini-0274:phyxminiproblems-problemphyxmini0274-energyinmillijoules, def:physics:phyx-mini-0274:phyxminiproblems-problemphyxmini0274-milliradianstoradians, def:physics:phyx-mini-0274:phyxminiproblems-problemphyxmini0274-simplependulumsetup}
  Numerical and qualitative evidence stated in the problem or read from the primary image. The image, rather than its auxiliary caption, places K\_s on the second horizontal grid line and the apex on the third, giving the 3 / 2 relation below. No pendulum-length value occurs here
\end{definition}
\begin{proof}
  This is the declared model definition of Matches Problem And Primary Figure.
\end{proof}
\begin{definition}[Uses Standard Earth Gravity]
  \label{def:physics:phyx-mini-0274:phyxminiproblems-problemphyxmini0274-usesstandardearthgravity}
  \lean{PhyXMiniProblems.ProblemPhyXMini0274.UsesStandardEarthGravity}
  \uses{def:physics:phyx-mini-0274:phyxminiproblems-problemphyxmini0274-accelerationinmeterspersecondsquared, def:physics:phyx-mini-0274:phyxminiproblems-problemphyxmini0274-simplependulumsetup}
  The standard near-Earth gravitational acceleration used by the answer
\end{definition}
\begin{proof}
  This is the declared model definition of Uses Standard Earth Gravity.
\end{proof}
\begin{definition}[Has Physical Pendulum Parameters]
  \label{def:physics:phyx-mini-0274:phyxminiproblems-problemphyxmini0274-hasphysicalpendulumparameters}
  \lean{PhyXMiniProblems.ProblemPhyXMini0274.HasPhysicalPendulumParameters}
  \uses{def:physics:phyx-mini-0274:phyxminiproblems-problemphyxmini0274-massinkilograms, def:physics:phyx-mini-0274:phyxminiproblems-problemphyxmini0274-lengthinmeters, def:physics:phyx-mini-0274:phyxminiproblems-problemphyxmini0274-accelerationinmeterspersecondsquared, def:physics:phyx-mini-0274:phyxminiproblems-problemphyxmini0274-energyinjoules, def:physics:phyx-mini-0274:phyxminiproblems-problemphyxmini0274-milliradianstoradians, def:physics:phyx-mini-0274:phyxminiproblems-problemphyxmini0274-simplependulumsetup}
  Positivity and nondegeneracy conditions for the physical experiment
\end{definition}
\begin{proof}
  This is the declared model definition of Has Physical Pendulum Parameters.
\end{proof}
\begin{definition}[Satisfies Ideal Pendulum Energy Law]
  \label{def:physics:phyx-mini-0274:phyxminiproblems-problemphyxmini0274-satisfiesidealpendulumenergylaw}
  \lean{PhyXMiniProblems.ProblemPhyXMini0274.SatisfiesIdealPendulumEnergyLaw}
  \uses{def:physics:phyx-mini-0274:phyxminiproblems-problemphyxmini0274-massinkilograms, def:physics:phyx-mini-0274:phyxminiproblems-problemphyxmini0274-lengthinmeters, def:physics:phyx-mini-0274:phyxminiproblems-problemphyxmini0274-accelerationinmeterspersecondsquared, def:physics:phyx-mini-0274:phyxminiproblems-problemphyxmini0274-energyinjoules, def:physics:phyx-mini-0274:phyxminiproblems-problemphyxmini0274-milliradianstoradians, def:physics:phyx-mini-0274:phyxminiproblems-problemphyxmini0274-simplependulumsetup}
  The exact ideal-simple-pendulum energy law on the interval shown in the figure. Conservation of mechanical energy and U(θ) - U(0) = m g L (1 - cos θ) give the decrease of kinetic energy away from the vertical. The plotted arc is visually parabola-like over its small angular range, but no small-angle approximation is promoted to an equality. This is a governing law, not the requested numerical length
\end{definition}
\begin{proof}
  This is the declared model definition of Satisfies Ideal Pendulum Energy Law.
\end{proof}
\begin{definition}[Answer Choice]
  \label{def:physics:phyx-mini-0274:phyxminiproblems-problemphyxmini0274-answerchoice}
  \lean{PhyXMiniProblems.ProblemPhyXMini0274.AnswerChoice}
  Labels printed beside the four candidate lengths
\end{definition}
\begin{proof}
  This is the declared model definition of Answer Choice.
\end{proof}
\begin{definition}[length In Meters]
  \label{def:physics:phyx-mini-0274:phyxminiproblems-problemphyxmini0274-answerchoice-lengthinmeters}
  \lean{PhyXMiniProblems.ProblemPhyXMini0274.AnswerChoice.lengthInMeters}
  \uses{def:physics:phyx-mini-0274:phyxminiproblems-problemphyxmini0274-lengthinmeters, def:physics:phyx-mini-0274:phyxminiproblems-problemphyxmini0274-answerchoice}
  Length in metres printed beside each answer label
\end{definition}
\begin{proof}
  This is the declared model definition of length In Meters.
\end{proof}
\begin{definition}[recorded Answer Choice]
  \label{def:physics:phyx-mini-0274:phyxminiproblems-problemphyxmini0274-recordedanswerchoice}
  \lean{PhyXMiniProblems.ProblemPhyXMini0274.recordedAnswerChoice}
  \uses{def:physics:phyx-mini-0274:phyxminiproblems-problemphyxmini0274-answerchoice}
  The answer label recorded in the supplied dataset metadata
\end{definition}
\begin{proof}
  This is the declared model definition of recorded Answer Choice.
\end{proof}
\begin{definition}[Matches Displayed Length]
  \label{def:physics:phyx-mini-0274:phyxminiproblems-problemphyxmini0274-matchesdisplayedlength}
  \lean{PhyXMiniProblems.ProblemPhyXMini0274.MatchesDisplayedLength}
  \uses{def:physics:phyx-mini-0274:phyxminiproblems-problemphyxmini0274-lengthinmeters, def:physics:phyx-mini-0274:phyxminiproblems-problemphyxmini0274-simplependulumsetup, def:physics:phyx-mini-0274:phyxminiproblems-problemphyxmini0274-answerchoice}
  Agreement with a displayed two-decimal length to the nearest centimetre
\end{definition}
\begin{proof}
  This is the declared model definition of Matches Displayed Length.
\end{proof}
\begin{definition}[Is Unique Closest Displayed Choice]
  \label{def:physics:phyx-mini-0274:phyxminiproblems-problemphyxmini0274-isuniqueclosestdisplayedchoice}
  \lean{PhyXMiniProblems.ProblemPhyXMini0274.IsUniqueClosestDisplayedChoice}
  \uses{def:physics:phyx-mini-0274:phyxminiproblems-problemphyxmini0274-lengthinmeters, def:physics:phyx-mini-0274:phyxminiproblems-problemphyxmini0274-simplependulumsetup, def:physics:phyx-mini-0274:phyxminiproblems-problemphyxmini0274-answerchoice}
  The selected length is strictly closer than every other displayed choice
\end{definition}
\begin{proof}
  This is the declared model definition of Is Unique Closest Displayed Choice.
\end{proof}
% --- Archon declaration topology end ---
% --- Archon physics formalization source end ---
